\chapter{Glossaire pratique}
\begin{description}
\item[RAW] Données de prise de vue à développer ; leur rendu dépend notamment
du décodage, de la balance des blancs et des profils.
\item[DCP] Profil couleur appareil utilisé dans le développement RAW.
\item[LCP] Profil objectif permettant certaines corrections optiques.
\item[ICC] Description colorimétrique d’un espace ou d’un périphérique ;
distincte d’une palette de couleurs de l’interface.
\item[EV ou IL] Unité de variation d’exposition ; une augmentation de 1 EV
correspond à un gain lumineux linéaire de facteur deux avant le rendu de sortie.
\item[Température] Réglage de balance des blancs exprimé en kelvins, à
distinguer d’une mesure universelle de tous les pixels de l’image.
\item[Nuance] Correction suivant l’axe vert/magenta.
\item[Teinte] Famille chromatique, par exemple rouge ou orange ; ne pas
confondre avec le nom du curseur Nuance.
\item[Luminance] Grandeur relative décrivant le niveau lumineux ; sa définition
précise dépend de l’espace et de l’étape de calcul.
\item[Gamut] Ensemble des couleurs représentables dans un espace donné.
\item[sRGB] Espace de sortie courant des exports rendus décrits dans ce livre.
\item[Aperçu] Image calculée pour la visualisation, éventuellement moins définie
que la source ; ce n’est pas la preuve de résolution de l’export.
\item[Compagnon ou sidecar] Fichier associé à la source, ici \chemin{.ipr},
conservant des réglages sans remplacer les pixels du fichier source.
\item[INI] Fichier des préférences d’interface, indépendant des compagnons photo.
\item[Neutre] Référence de réglages déclarée pour une comparaison ; elle doit
toujours préciser le développement de départ, le profil et la balance des blancs.
\end{description}
